\section{Killing spinors and their superalgebras}\label{sec:spin}
For either non-null choice of the AdS fibre, every regular deformation with
$\gamma\ne0$ has exactly eight real Killing spinors. For the null choice
there are at least four, and exactly four on an open dense set of parameters.
We work locally, or on the simply connected cover; descent to a quotient
also requires invariance under the lifted identifications.

\subsection{Two mixed curvatures for a non-null fibre}
Use the homothetic fields $(\hat g,\hat F)$ in \eqref{eq:kse}, and set
$k=\sqrt G$, $b=s\gamma k$, so that $k^2+s b^2=1$ and $k>0$.
In each block $i=0,+,-$, use an oriented seed group coframe
$(h_i^1,h_i^2,\theta_i)$, with the last axis along the selected fibre.
In the hatted metric the horizontal coframes are unchanged and the unit
fibre coframe is $k\theta_i$. Write $f_i$ for its index, $e_a$ for the
hatted frame vectors and $\sigma_a=\hat g(e_a,e_a)$. Choose group frames
with seed Maurer--Cartan equations $d e_{(0)}^c=\sigma_c q_i
 e_{(0)}^a\wedge e_{(0)}^d$ for cyclic axes $(a,d,c)$ in a block.
Then $[e_a,e_d]=C_{ad}{}^c e_c$, where for those cyclic axes
\[
 C_{ad}{}^c=-q_i\sigma_c m_c,\qquad m_h=k^{-1},\quad m_f=k .
\]
Brackets between blocks vanish. With upper-index gamma matrices satisfying
$\{\Gamma^a,\Gamma^d\}=2\hat g^{ad}$ in this frame, put
\[
 \begin{gathered}
 \omega_i=\Gamma^{h_i^1h_i^2f_i},\qquad N=\sum_i q_i\omega_i,\qquad K=\Gamma^{f_0f_+f_-},\\
 X=\Gamma^x,\qquad Y=\Gamma^y,\qquad Z_+=X+bK,\qquad Z_-=-X+bK .
 \end{gathered}
\]
In \eqref{eq:fields} the seed flux contributes $NX$ to
$\not{\hat F}$, and the three terms differentiating the fibre volume
contribute $bNK$: $\omega_i K$ removes $f_i$ with its metric sign
and the sign of its position in the ordered triple. Thus
$\not{\hat F}=NZ_+$. Distinct $\omega_i$ anticommute, and their squares
are $+1,-1,-1$, respectively. Hence
\[
 \begin{gathered}
 N^2=0,\qquad [N,K]=0,\qquad \{N,X\}=\{N,Y\}=0,\\
 K^2=-s,\qquad Z_\pm^2=k^2,\qquad Z_\pm N=NZ_\mp .
 \end{gathered}
\]
Since $\{N,\omega_0\}=2q_0$, every element of $\ker N$ lies in
$\operatorname{im}N$; the two coincide and have real dimension sixteen.

For an axis in block $i$ write $J_a=\Gamma_a\omega_i$, so
$J_a=\Gamma^d\Gamma^c$ for cyclic $(a,d,c)$. The Koszul formula gives
$2\langle\nabla_a e_d,e_c\rangle=-q_i(m_c-m_a+m_d)$, and hence the spin
connection matrix is $q_i(m_d+m_c-m_a)J_a/4$. The flux matrix in
\eqref{eq:kse} is $-\Gamma_a\not{\hat F}/24+\not{\hat F}\Gamma_a/8$.
Adding them, write $D_a=e_a+\mathcal A_a$, with
\begin{equation}\label{eq:spin-connection}
\begin{aligned}
 \mathcal A_h&=\tfrac1{12}\Gamma_h NZ_+
                 +\tfrac{q_i}{4}J_h(k-Z_+),\\
 \mathcal A_f&=\tfrac1{12}\Gamma_f N(X-2bK)
                 +\tfrac{q_i}{4}J_f(2/k-k+Z_-),\\
 \mathcal A_y&=\tfrac1{12}YNZ_+,\qquad
 \mathcal A_x=\tfrac1{12}N(2-bXK).
\end{aligned}
\end{equation}
All $J_a$ commute with $N$ and $X$; $J_h$ anticommutes with $K$,
whereas $J_f$ commutes with it. The matrices are constant, so the curvature
on an arbitrary spinor is $\mathcal R_{ad}=[\mathcal A_a,\mathcal A_d]
-C_{ad}{}^c\mathcal A_c$. In particular, no assumption on coordinate
dependence enters the following two necessary conditions.

First take $f=f_i$ and the flat direction $y$. There is no frame bracket,
and the Levi--Civita term in $\mathcal A_f$ commutes with $\mathcal A_y$.
Set $L_f=X-2bK$. Using $N\Gamma_fN=2q_iJ_f N$ and $Z_+N=NZ_-$,
the two remaining commutators in \eqref{eq:spin-connection} give
\begin{equation}\label{eq:spin-fy}
 \mathcal R_{fy}=-\frac{q_i}{144}YJ_f N
 \big[2Z_+L_f+3(Z_+^2+Z_+Z_-)\big]
 =-\frac{q_i k^2}{72}YJ_f N .
\end{equation}
Indeed $2L_f+3(Z_++Z_-)=2Z_+$. As $YJ_f$ is invertible, every solution
lies in $\ker N$. This constant subbundle is preserved by the connection:
$X,K,Z_\pm$ preserve it, so the terms with left factor $\Gamma_a N$ and
both flat matrices vanish there.

Next take $f_i$ and $h_j$ in distinct blocks. Their frame bracket vanishes,
$J_{f_i}$ commutes with $J_{h_j}$, and $Z_-J_{h_j}=-J_{h_j}Z_+$.
With $d=2/k-k$ the restricted curvature is
\begin{equation}\label{eq:spin-fh}
\begin{aligned}
 \mathcal R_{f_i h_j}\big|_{\ker N}
 &=\frac{q_iq_j}{16}J_{f_i}J_{h_j}
   \big[(d-Z_+)(k-Z_+)-(k-Z_+)(d+Z_-)\big]\\
 &=-\frac{bq_iq_j}{8}J_{f_i}J_{h_j}(k-Z_+)K .
\end{aligned}
\end{equation}
For $b\ne0$, conjugation by $K$ interchanges $Z_+$ and $Z_-$, so the
allowed initial values lie in $E=\ker N\cap\ker(Z_--k)$. On $\ker N$,
$Z_-^2=k^2$ and the invertible $Y$ exchanges the two real eigenspaces of
$Z_-$. Thus $\dim E=8$, including throughout the regular timelike interval.

\subsection{The eight solutions and their currents}
At $b=0$ let $E_0=\ker N\cap\ker(X+1)$. Its seed connection matrices
are $q_iJ_a/2$ on each group and zero on the plane. This connection is
flat: for cyclic axes in block $i$, $[J_a,J_d]=-2\sigma_c J_c$, and
matrices from distinct blocks commute. Every initial value in $E_0$
therefore gives a parallel spinor. The complementary $X=+1$ eigenspace
in $\ker N$ has zero connection matrices, giving sixteen seed spinors in
all. To transport the first eight to $b\ne0$, set
\begin{equation}\label{eq:spin-intertwiner}
 u=\frac{b}{1+k},\qquad S=\sqrt{\frac{1+k}{2}}\,(1+uK).
\end{equation}
The identity $1+su^2=2/(1+k)$ makes $S$ invertible. It commutes with $N$;
on $X=-1$ one has $Z_-(1+uK)=(1-sbu)+(b-u)K=k(1+uK)$.
Thus $S E_0=E$. On $E$, substituting $Z_+=2bK-Z_-$ into
\eqref{eq:spin-connection} gives
\[
 \mathcal A_h=\tfrac{q_i}{2}J_h(k-bK),\qquad
 \mathcal A_f=\tfrac{q_i}{2k}J_f,\qquad \mathcal A_x=\mathcal A_y=0 .
\]
Now $(k-bK)(1+uK)=1-uK$, so $J_h(k-bK)S=SJ_h$, while $J_fS=SJ_f$.
The hatted fibre vector is $e_f=e_f^{(0)}/k$; multiplying its differential
equation by $k$ gives the seed equation on the same vector field.
Consequently $\epsilon=S\epsilon_0$, with components compared in the two
group frames, constructs all eight solutions.

To compute their polarized Dirac currents, use the Majorana pairing $B$
for which one-blades are skew-adjoint and three-blades are self-adjoint.
In particular $K$ is self-adjoint. For a horizontal axis $h$ and a fibre
axis $f$, the matrix between the two seed spinors in the current is
\[
 S\Gamma^h S=\Gamma^h,\qquad
 S\Gamma^f S=\tfrac{1+k}{2}\Gamma^f(1-su^2+2uK).
\]
The term $\Gamma^f K$ commutes with $X$ and hence pairs two $X=-1$ spinors
to zero by skew-adjointness of $X$. Since $1-su^2=k(1+su^2)$, the fibre
component acquires a factor $k$, cancelled in the vector current by
$e_f=e_f^{(0)}/k$. Moreover $\ker N=\operatorname{im}N$ is totally
isotropic for $B$, since $N$ is self-adjoint and $N^2=0$. Both $X$ and
$Y$ preserve this kernel, so there is no flat component. Thus the currents
in the hatted metric are exactly the seed vector fields, not just vectors
with the same components in two different frames. Restoring $g$ changes
only their common normalization.

\subsection{The four invariant spinors in the null branch}
Call the seed symmetry sectors $L,R$ and choose $H_+,H_-$ in $R$ and the
parabolic generator $n$ in $L$. The odd module of each sector is the real
form of $(\mathbf 2,\mathbf 2,\mathbf 2)$ for its three simple factors.
Each selected Hopf has no zero weight on its sphere doublet, leaving no
$R$ spinor invariant under both. In $L$ the sphere conditions are trivial,
and the spinorial Lie derivative along $n$ is rank-one nilpotent on the
AdS doublet. Its kernel times the two sphere doublets gives four real
spinors invariant under $W,H_+,H_-$, since the seed spinors are independent
of $y$.

Explicitly, use the null-branch seed frame with blocks $012,345,678$,
$0$ timelike, $x=9$, $y=10$, $n=-e_0+e_1$ and $H_+=e_5$, $H_-=e_8$.
The projector onto these four solutions at a point is
\begin{equation}\label{eq:spin-null-projector}
 P_0=\tfrac18(1+\Gamma^9)(1+\Gamma^{01})
 \left[1-\frac{\Gamma^{012}(q_+\Gamma^{345}+q_-\Gamma^{678})}{q_0}\right].
\end{equation}
The three commuting factors select the sphere-invariant handedness, the
line $\not p\,\epsilon_0=0$, and $\ker N$, respectively. For $c\ne0$ the
three reduction directions span a spacelike plane. In the hatted coframe
obtained by scaling that plane by $D^{-1/2}$, consider the following
spinor map (an overall constant does not affect the Killing-spinor equation):
\begin{equation}\label{eq:spin-null-rotation}
 S_n=1+\frac{\gamma}{1+\sqrt D}\,\not w\,\not\eta_+\,\not\eta_- ,
\end{equation}
where the slashes use the seed frame. The triple squares to $-c^2$, so for $c\ne0$ the normalized matrix is a rotation. Its coefficient and sign are fixed by the connection identity below. The coframe
extends through $c=0$, since $(D^{-1/2}-1)/c^2=-\gamma^2/[\sqrt D(1+
\sqrt D)]$, and $S_n$ stays invertible there. In the homogeneous frames,
write $e'_a=N^b{}_a e_b^{(0)}$ with $N^b{}_a$ the inverse coframe matrix of Appendix~\ref{app:certificate} and let $\mathcal A'_a,\mathcal A_b^{(0)}$
be the corresponding connection matrices. The differential statement is
\[
 (\mathcal A'_a S_n-S_n N^b{}_a\mathcal A_b^{(0)})P_0=0
 \qquad (a=0,\ldots,10).
\]
Appendix~\ref{app:certificate} proves these identities for the family;
hence $S_n\epsilon_0$ gives four independent solutions, including at $c=0$.
The same appendix exhibits a $28\times28$ minor of the stacked curvature
matrices which is nonzero at one regular point. Its entries are analytic
on the connected parameter domain, so the common Majorana kernel has
dimension at most four on an open dense set. This proves generic equality,
not enhancement wherever the chosen minor vanishes.

\subsection{Odd brackets and spectator isometries}
We distinguish the algebra generated by the odd part from the extension by
all flux-preserving Killing vectors, using the Dirac current and spinorial
Lie derivative as brackets~\cite{FOF}; no three-form gauge generators are
included. To fix the parameter of the seed exceptional algebra, put
$L_i=\omega_0\omega_i$ for $i=+,-$. These are anticommuting involutions,
and the projector onto $\ker N$ is $P_*=(1-\sum_{i=+,-}q_i L_i/q_0)/2$.
Thus $P_*L_iP_*=-(q_i/q_0)P_*$. This compression can be used inside the
current: $P_*$ is self-adjoint for $B(\cdot,\omega_0\cdot)$, and $J_a$
commutes with it. In AdS one has $\Gamma_a=\omega_0J_a$; on sphere $i$,
$\Gamma_a=-\omega_0 L_i J_a$. The sphere current therefore carries a
factor $q_i/q_0$ relative to the AdS current. Expressing the unit vectors
in generators with radius-independent brackets supplies a second factor
$L/R_i=q_i/q_0$. In the convention with the noncompact coefficient
negative, the three bracket coefficients have ratios
\begin{equation}\label{eq:spin-alpha}
 -q_0^2:q_+^2:q_-^2=-(1+\alpha):1:\alpha,\qquad
 \alpha=\frac{q_-^2}{q_+^2}=\frac{R_+^2}{R_-^2} .
\end{equation}
The sum vanishes as required by the odd Jacobi identity. Exchanging sphere
labels inverts $\alpha$.

The seed has two commuting $\mathfrak d(2,1;\alpha)$ sectors~\cite{GGP}. For a
nonzero non-null deformation the intertwiner \eqref{eq:spin-intertwiner}
is constant in group frames and equivariant under the three surviving
translation groups. Its current calculation retains their seed brackets
with the same $\alpha$. The odd-generated algebra is consequently one
$\mathfrak d(2,1;\alpha)$, whereas the full local Killing superalgebra is
\[
 \mathfrak d(2,1;\alpha)\oplus\mathbb R^5 .
\]
Within each non-round block the Ricci tensor distinguishes the fibre line; a connected isometry therefore preserves the fibre metric and the round metric reconstructed from it. On the opposite side of the translation group this leaves only the selected one-parameter subgroup. Nor can a connected isometry mix different curved blocks, even for equal radii. By the three-dimensional curvature identity, each horizontal--fibre plane has nonzero sectional curvature $\delta_i a_i^2/4$ in the notation of Section~\ref{sec:constraints}. The two corresponding curvature endomorphisms generate the block rotation algebra, so each curved tangent block is an irreducible holonomy summand, separate from the flat directions. Equal blocks allow at most a discrete interchange. These three fibre actions commute with the rotation $S$
(they preserve the ordered fibre volume) and act trivially on the selected
seed sector. The two flat translations also act trivially; a rotation of
the plane does not preserve the $V_0\wedge dx$ flux. The abelian summands
here and below denote Lie algebras, not global identifications.

At a generic null point with $c\gamma\ne0$, the four spinors transform as
the real $(\mathbf 2,\mathbf 2)$ of the two full sphere groups. They still
obey $\not p\,\epsilon=0$: $n$ is orthogonal to the deformation plane,
and the triple in \eqref{eq:spin-null-rotation} anticommutes with
$\not p$. The Dirac current of a nonzero Majorana spinor is nonzero and
causal. Being orthogonal to $n$, it is therefore proportional to $n$;
its coefficient is constant since $fn$ is Killing only if
$df\otimes p+p\otimes df=0$. Polarization and the unique invariant
symmetric inner product of $(\mathbf 2,\mathbf 2)$ give, in a normalized
basis,
\[
 \{Q_I,Q_J\}=\delta_{IJ} n,\qquad [n,Q_I]=0,\qquad I,J=1,\ldots,4 .
\]
The odd-generated algebra is the supertranslation algebra
$\mathfrak t^{1|4}$, with even line $\mathbb R n$, not $\mathbb R W$.
For the full algebra, the non-Einstein part of the null AdS Ricci tensor
fixes $p\otimes p$. Thus a connected isometry of this block preserves
$p$, hence also the seed AdS metric reconstructed by adding back the wave
term. The stabilizer keeps the opposite AdS group and the line generated
by $n$. Each Berger sphere keeps its translation group and its selected
Hopf subgroup; no curved block has a parallel direction that can mix
with the flat plane. The selected Hopfs, opposite AdS group and flat
translations preserve the frame transformation and \eqref{eq:spin-null-rotation}
and act trivially on the four invariant seed spinors. The full algebra is
therefore
\begin{equation}\label{eq:spin-null-algebra}
 \big[(\mathfrak{su}(2)_{+,L}\oplus\mathfrak{su}(2)_{-,L})
       \ltimes\mathfrak t^{1|4}\big]
 \oplus\mathfrak{sl}(2,\mathbb R)_R\oplus\mathbb R^4 ,
\end{equation}
where the last four generators are the two selected Hopfs and the two flat
translations in the unsheared coordinates of \eqref{eq:fields}.
